\subsection{Définition structurelle et fonction de la constante gravitationnelle}
\label{sec:ii-definicion-gravedad}

\begin{definition}[Section gravitationnelle contragrédiente de l'action]
Pour une section positive \(\hbar_a\), une vitesse
\(c_{\rm int}>0\) et le rayon chronologique \(R_{108}\) de la
section~\ref{sec:gravity}, la section gravitationnelle circulaire est
\begin{equation}
 G_a=\frac{c_{\rm int}^{3}R_{108}^{2}}{\hbar_a},
 \qquad R_{108}=\frac{54}{\pi}\ell_0,
 \qquad \ell_0=c_{\rm int}t_0.
 \label{eq:ii-definicion-gravedad-seccion}
\end{equation}
C'est la section réciproque de l'action qui réalise la coïncidence
entre le rayon gravitationnel réduit du cycle et son rayon
chronologique, dans la convention \(r_g=Gm/c_{\rm int}^2\).
\end{definition}

La proposition de coïncidence des rayons de la
section~\ref{sec:gravity} démontre son unicité dans la collection
\(G_a(\vartheta)=\vartheta c_{\rm int}^{5}t_0^2/\hbar_a\).
La condition circulaire fixe \(\vartheta=(54/\pi)^2\) et conduit
à \eqref{eq:ii-definicion-gravedad-seccion}. Cette condition et les
sections dimensionnelles positives font partie de la réalisation.
La section~\ref{md:rigidez} construit sa longueur et réunit la
preuve d’unicité pour le caractère positif général, avec
$R_{108}=L_\alpha$. Ainsi le sélecteur circulaire et le lecteur
longitudinal conservent une même normalisation.
En les maintenant communes aux deux orientations, on obtient
\begin{equation}
 \frac{G_{\mathrm{pre}}}{G_{\mathrm{ret}}}
 =R_{\mathrm{act}}^{-1},
 \qquad
 \frac{G_a\hbar_a}{c_{\rm int}^{3}}
 =R_{108}^{2}=\ell_P^2.
 \label{eq:ii-definicion-gravedad-area}
\end{equation}
La preuve est la substitution directe de
\eqref{eq:ii-definicion-gravedad-seccion} ; le changement d'action
est compensé par le changement réciproque de \(G_a\).
